\documentclass[12pt]{article}

\usepackage{float}
\usepackage{graphicx}			% Use this package to include images
\usepackage{amsmath}			% A library of many standard math expressions
\usepackage[margin=1in]{geometry}% Sets 1in margins. 
\usepackage{fancyhdr}			% Creates headers and footers
\usepackage{enumerate}          %These two package give custom labels to a list
\usepackage[shortlabels]{enumitem}
\usepackage{matlab-prettifier}
\usepackage{tikz}
\usepackage[utf8]{inputenc}
\usepackage{pgfplots} 
\usepackage{pgfgantt}
\usepackage{pdflscape}
\usepackage{array,multirow}
\usepackage{tabularx,booktabs}
\newcolumntype{Y}{>{\centering\arraybackslash}X}
\usetikzlibrary{shapes.geometric}
\usetikzlibrary{arrows.meta,arrows}
\pgfplotsset{compat=newest} 
\pgfplotsset{plot coordinates/math parser=false}
\newlength\fwidth
\newlength\fheight
\usepackage{float,graphicx}
\pagestyle{fancy}
\fancyhead[l]{}
\fancyhead[c]{Experiment 7}
\fancyhead[r]{}
\fancyfoot[c]{\thepage}
\setlength{\headheight}{20pt}
\begin{document}
\author{}
\date{}
\title{Brewster's Angle}
\maketitle
\section*{Aim}
To find refractive index of glass by measuring brewster angle.
\section*{Analysis}
For the first reading, the vertical alignment changed linearly as the angle was changing which means that the amount of light measured needs to be scaled up as it deviated from the minimum value.
We can reasonably approximate the distribution of intensity of a laser by a gaussian.
Fitting a gaussian for the data we get the following plot.
\begin{figure}[H]
    \centering
    \setlength\fwidth{0.75\textwidth}
    \setlength\fheight{0.5\textwidth}
    % This file was created by matlab2tikz.
%
\definecolor{mycolor1}{rgb}{0.00000,0.44700,0.74100}%
\definecolor{mycolor2}{rgb}{0.12941,0.12941,0.12941}%
%
\begin{tikzpicture}

\begin{axis}[%
width=0.951\fwidth,
height=\fheight,
at={(0\fwidth,0\fheight)},
scale only axis,
xmin=-50,
xmax=50,
xlabel style={font=\color{mycolor2}},
xlabel={Angle (in degrees)},
ymin=0,
ymax=20,
ylabel style={font=\color{mycolor2}},
ylabel={Intensity (in mA)},
axis background/.style={fill=white},
title style={font=\bfseries\color{mycolor2}},
title={$\text{Distribution (}\sigma\text{ = 10.884)}$},
xmajorgrids,
ymajorgrids,
legend style={legend cell align=left, align=left}
]
\addplot [color=mycolor1]
  table[row sep=crcr]{%
-50	0.000460078048947578\\
-49.8998998998999	0.000479912674243277\\
-49.7997997997998	0.000500560056954716\\
-49.6996996996997	0.000522051596599788\\
-49.5995995995996	0.000544419823420536\\
-49.4994994994995	0.000567698435689168\\
-49.3993993993994	0.000591922338109637\\
-49.2992992992993	0.000617127681341812\\
-49.1991991991992	0.000643351902675775\\
-49.0990990990991	0.000670633767884217\\
-48.998998998999	0.000699013414281412\\
-48.8988988988989	0.000728532395017747\\
-48.7987987987988	0.000759233724639237\\
-48.6986986986987	0.000791161925941933\\
-48.5985985985986	0.000824363078151659\\
-48.4984984984985	0.00085888486645993\\
-48.3983983983984	0.000894776632947465\\
-48.2982982982983	0.000932089428927083\\
-48.1981981981982	0.000970876068738363\\
-48.0980980980981	0.0010111911850268\\
-47.997997997998	0.00105309128554084\\
-47.8978978978979	0.00109663481148037\\
-47.7977977977978	0.00114188219743101\\
-47.6976976976977	0.00118889593291883\\
-47.5975975975976	0.00123774062562056\\
-47.4974974974975	0.00128848306626492\\
-47.3973973973974	0.00134119229526108\\
-47.2972972972973	0.00139593967109069\\
-47.1971971971972	0.00145279894050035\\
-47.0970970970971	0.00151184631053185\\
-46.996996996997	0.00157316052242787\\
-46.8968968968969	0.00163682292745133\\
-46.7967967967968	0.00170291756465684\\
-46.6966966966967	0.00177153124065318\\
-46.5965965965966	0.00184275361139605\\
-46.4964964964965	0.00191667726605093\\
-46.3963963963964	0.00199339781296561\\
-46.2962962962963	0.00207301396779331\\
-46.1961961961962	0.00215562764380639\\
-46.0960960960961	0.00224134404444221\\
-45.995995995996	0.00233027175812187\\
-45.8958958958959	0.00242252285538374\\
-45.7957957957958	0.0025182129883732\\
-45.6956956956957	0.002617461492731\\
-45.5955955955956	0.00272039149192201\\
-45.4954954954955	0.00282713000404714\\
-45.3953953953954	0.00293780805118084\\
-45.2952952952953	0.00305256077127702\\
-45.1951951951952	0.00317152753268615\\
-45.0950950950951	0.00329485205132659\\
-44.994994994995	0.00342268251055332\\
-44.8948948948949	0.0035551716837669\\
-44.7947947947948	0.00369247705980604\\
-44.6946946946947	0.00383476097116684\\
-44.5945945945946	0.00398219072509176\\
-44.4944944944945	0.00413493873757133\\
-44.3943943943944	0.00429318267030165\\
-44.2942942942943	0.0044571055706403\\
-44.1941941941942	0.00462689601460352\\
-44.0940940940941	0.00480274825294689\\
-43.993993993994	0.00498486236037177\\
-43.8938938938939	0.00517344438789945\\
-43.7937937937938	0.00536870651845469\\
-43.6936936936937	0.00557086722569968\\
-43.5935935935936	0.00578015143615947\\
-43.4934934934935	0.00599679069467906\\
-43.3933933933934	0.00622102333325249\\
-43.2932932932933	0.00645309464326277\\
-43.1931931931932	0.00669325705117179\\
-43.0930930930931	0.00694177029769847\\
-42.992992992993	0.00719890162052228\\
-42.8928928928929	0.00746492594054935\\
-42.7927927927928	0.00774012605177678\\
-42.6926926926927	0.00802479281479058\\
-42.5925925925926	0.00831922535393144\\
-42.4924924924925	0.00862373125816143\\
-42.3923923923924	0.0089386267856644\\
-42.2922922922923	0.00926423707221086\\
-42.1921921921922	0.00960089634331778\\
-42.0920920920921	0.00994894813023214\\
-41.991991991992	0.010308745489766\\
-41.8918918918919	0.0106806512280092\\
-41.7917917917918	0.0110650381279454\\
-41.6916916916917	0.0114622891809937\\
-41.5915915915916	0.0118727978224996\\
-41.4914914914915	0.0122969681711943\\
-41.3913913913914	0.0127352152726418\\
-41.2912912912913	0.0131879653466907\\
-41.1911911911912	0.0136556560389466\\
-41.0910910910911	0.0141387366762767\\
-40.990990990991	0.0146376685263606\\
-40.8908908908909	0.0151529250612943\\
-40.7907907907908	0.0156849922252566\\
-40.6906906906907	0.0162343687062415\\
-40.5905905905906	0.0168015662118602\\
-40.4904904904905	0.0173871097492134\\
-40.3903903903904	0.0179915379088314\\
-40.2902902902903	0.0186154031526777\\
-40.1901901901902	0.0192592721062093\\
-40.0900900900901	0.0199237258544833\\
-39.98998998999	0.020609360242298\\
-39.8898898898899	0.0213167861783529\\
-39.7897897897898	0.0220466299434088\\
-39.6896896896897	0.0227995335024272\\
-39.5895895895896	0.0235761548206637\\
-39.4894894894895	0.0243771681836888\\
-39.3893893893894	0.0252032645213043\\
-39.2892892892893	0.0260551517353202\\
-39.1891891891892	0.0269335550311556\\
-39.0890890890891	0.0278392172532215\\
-38.988988988989	0.0287728992240389\\
-38.8888888888889	0.0297353800870457\\
-38.7887887887888	0.0307274576530372\\
-38.6886886886887	0.0317499487501851\\
-38.5885885885886	0.0328036895775726\\
-38.4884884884885	0.0338895360621817\\
-38.3883883883884	0.0350083642192628\\
-38.2882882882883	0.0361610705160136\\
-38.1881881881882	0.0373485722384879\\
-38.0880880880881	0.0385718078616531\\
-37.987987987988	0.0398317374225089\\
-37.8878878878879	0.0411293428961744\\
-37.7877877877878	0.0424656285748482\\
-37.6876876876877	0.0438416214495384\\
-37.5875875875876	0.0452583715944576\\
-37.4874874874875	0.0467169525539695\\
-37.3873873873874	0.0482184617319718\\
-37.2872872872873	0.0497640207835921\\
-37.1871871871872	0.0513547760090706\\
-37.0870870870871	0.0529918987496952\\
-36.986986986987	0.0546765857856524\\
-36.8868868868869	0.0564100597356489\\
-36.7867867867868	0.0581935694581557\\
-36.6866866866867	0.0600283904541162\\
-36.5865865865866	0.0619158252709626\\
-36.4864864864865	0.0638572039077676\\
-36.3863863863864	0.0658538842213639\\
-36.2862862862863	0.0679072523332475\\
-36.1861861861862	0.0700187230370848\\
-36.0860860860861	0.0721897402066271\\
-35.985985985986	0.0744217772038405\\
-35.8858858858859	0.0767163372870418\\
-35.7857857857858	0.0790749540188357\\
-35.6856856856857	0.0814991916736334\\
-35.5855855855856	0.0839906456445301\\
-35.4854854854855	0.0865509428493145\\
-35.3853853853854	0.089181742135371\\
-35.2852852852853	0.0918847346832355\\
-35.1851851851852	0.0946616444085526\\
-35.0850850850851	0.0975142283621819\\
-34.984984984985	0.100444277128188\\
-34.8848848848849	0.103453615219449\\
-34.7847847847848	0.106544101470602\\
-34.6846846846847	0.109717629428052\\
-34.5845845845846	0.112976127736744\\
-34.4844844844845	0.116321560523418\\
-34.3843843843844	0.119755927776021\\
-34.2842842842843	0.123281265718999\\
-34.1841841841842	0.126899647184121\\
-34.0840840840841	0.130613181976535\\
-33.983983983984	0.134424017235712\\
-33.8838838838839	0.138334337790961\\
-33.7837837837838	0.142346366511146\\
-33.6836836836837	0.146462364648285\\
-33.5835835835836	0.150684632174654\\
-33.4834834834835	0.155015508113041\\
-33.3833833833834	0.159457370859795\\
-33.2832832832833	0.164012638500266\\
-33.1831831831832	0.168683769116292\\
-33.0830830830831	0.173473261085314\\
-32.982982982983	0.178383653370747\\
-32.8828828828829	0.183417525803198\\
-32.7827827827828	0.188577499352136\\
-32.6826826826827	0.193866236387588\\
-32.5825825825826	0.199286440931466\\
-32.4824824824825	0.204840858898081\\
-32.3823823823824	0.21053227832344\\
-32.2822822822823	0.216363529582873\\
-32.1821821821822	0.222337485596574\\
-32.0820820820821	0.22845706202259\\
-31.981981981982	0.234725217436837\\
-31.8818818818819	0.241144953499672\\
-31.7817817817818	0.247719315108576\\
-31.6816816816817	0.254451390536484\\
-31.5815815815816	0.261344311555301\\
-31.4814814814815	0.268401253544132\\
-31.3813813813814	0.275625435581768\\
-31.2812812812813	0.283020120522929\\
-31.1811811811812	0.290588615057812\\
-31.0810810810811	0.298334269754454\\
-30.980980980981	0.306260479083419\\
-30.8808808808809	0.31437068142434\\
-30.7807807807808	0.322668359053812\\
-30.6806806806807	0.331157038114162\\
-30.5805805805806	0.339840288562594\\
-30.4804804804805	0.348721724100225\\
-30.3803803803804	0.357805002080515\\
-30.2802802802803	0.367093823396597\\
-30.1801801801802	0.376591932347021\\
-30.0800800800801	0.386303116479408\\
-29.97997997998	0.396231206411528\\
-29.8798798798799	0.406380075629314\\
-29.7797797797798	0.416753640261303\\
-29.6796796796797	0.427355858829039\\
-29.5795795795796	0.438190731972929\\
-29.4794794794795	0.449262302153084\\
-29.3793793793794	0.460574653324644\\
-29.2792792792793	0.472131910587122\\
-29.1791791791792	0.48393823980729\\
-29.0790790790791	0.495997847215118\\
-28.978978978979	0.508314978972322\\
-28.8788788788789	0.520893920713036\\
-28.7787787787788	0.533738997056165\\
-28.6786786786787	0.546854571088956\\
-28.5785785785786	0.560245043821347\\
-28.4784784784785	0.573914853610646\\
-28.3783783783784	0.587868475556115\\
-28.2782782782783	0.602110420863032\\
-28.1781781781782	0.616645236175803\\
-28.0780780780781	0.631477502879726\\
-27.977977977978	0.646611836371\\
-27.8778778778779	0.66205288529459\\
-27.7777777777778	0.677805330749564\\
-27.6776776776777	0.693873885461536\\
-27.5775775775776	0.710263292921835\\
-27.4774774774775	0.726978326493092\\
-27.3773773773774	0.744023788480854\\
-27.2772772772773	0.761404509170944\\
-27.1771771771772	0.779125345832233\\
-27.0770770770771	0.79719118168453\\
-26.976976976977	0.815606924831305\\
-26.8768768768769	0.834377507156975\\
-26.7767767767768	0.853507883188493\\
-26.6766766766767	0.873003028921016\\
-26.5765765765766	0.892867940607398\\
-26.4764764764765	0.913107633511332\\
-26.3763763763764	0.933727140623928\\
-26.2762762762763	0.954731511343568\\
-26.1761761761762	0.976125810118873\\
-26.0760760760761	0.99791511505465\\
-25.975975975976	1.0201045164807\\
-25.8758758758759	1.04269911548341\\
-25.7757757757758	1.06570402239998\\
-25.6756756756757	1.08912435527534\\
-25.5755755755756	1.11296523828165\\
-25.4754754754755	1.13723180010031\\
-25.3753753753754	1.16192917226666\\
-25.2752752752753	1.18706248747723\\
-25.1751751751752	1.2126368778597\\
-25.0750750750751	1.23865747320553\\
-24.974974974975	1.26512939916557\\
-24.8748748748749	1.29205777540849\\
-24.7747747747748	1.3194477137425\\
-24.6746746746747	1.34730431620018\\
-24.5745745745746	1.375632673087\\
-24.4744744744745	1.40443786099339\\
-24.3743743743744	1.43372494077098\\
-24.2742742742743	1.46349895547287\\
-24.1741741741742	1.49376492825872\\
-24.0740740740741	1.52452786026452\\
-23.973973973974	1.55579272843779\\
-23.8738738738739	1.58756448333834\\
-23.7737737737738	1.61984804690515\\
-23.6736736736737	1.65264831018979\\
-23.5735735735736	1.68597013105679\\
-23.4734734734735	1.71981833185162\\
-23.3733733733734	1.75419769703661\\
-23.2732732732733	1.78911297079558\\
-23.1731731731732	1.82456885460762\\
-23.0730730730731	1.8605700047908\\
-22.972972972973	1.89712103001628\\
-22.8728728728729	1.93422648879368\\
-22.7727727727728	1.97189088692837\\
-22.6726726726727	2.0101186749514\\
-22.5725725725726	2.04891424552287\\
-22.4724724724725	2.08828193080957\\
-22.3723723723724	2.12822599983775\\
-22.2722722722723	2.16875065582185\\
-22.1721721721722	2.2098600334701\\
-22.0720720720721	2.25155819626799\\
-21.971971971972	2.2938491337405\\
-21.8718718718719	2.33673675869416\\
-21.7717717717718	2.38022490443988\\
-21.6716716716717	2.42431732199776\\
-21.5715715715716	2.4690176772848\\
-21.4714714714715	2.51432954828676\\
-21.3713713713714	2.56025642221529\\
-21.2712712712713	2.60680169265145\\
-21.1711711711712	2.65396865667699\\
-21.0710710710711	2.70176051199441\\
-20.970970970971	2.75018035403731\\
-20.8708708708709	2.79923117307216\\
-20.7707707707708	2.84891585129284\\
-20.6706706706707	2.89923715990948\\
-20.5705705705706	2.95019775623268\\
-20.4704704704705	3.00180018075488\\
-20.3703703703704	3.05404685422997\\
-20.2702702702703	3.10694007475287\\
-20.1701701701702	3.16048201484052\\
-20.0700700700701	3.21467471851561\\
-19.96996996997	3.26952009839489\\
-19.8698698698699	3.32501993278338\\
-19.7697697697698	3.38117586277622\\
-19.6696696696697	3.43798938936967\\
-19.5695695695696	3.49546187058299\\
-19.4694694694695	3.55359451859281\\
-19.3693693693694	3.61238839688158\\
-19.2692692692693	3.67184441740198\\
-19.1691691691692	3.7319633377589\\
-19.0690690690691	3.79274575841064\\
-18.968968968969	3.85419211989127\\
-18.8688688688689	3.91630270005573\\
-18.7687687687688	3.97907761134962\\
-18.6686686686687	4.04251679810529\\
-18.5685685685686	4.10662003386626\\
-18.4684684684685	4.17138691874151\\
-18.3683683683684	4.23681687679183\\
-18.2682682682683	4.30290915344969\\
-18.1681681681682	4.36966281297475\\
-18.0680680680681	4.43707673594682\\
-17.967967967968	4.50514961679799\\
-17.8678678678679	4.57387996138604\\
-17.7677677677678	4.64326608461076\\
-17.6676676676677	4.71330610807525\\
-17.5675675675676	4.78399795779399\\
-17.4674674674675	4.85533936194955\\
-17.3673673673674	4.92732784869987\\
-17.2672672672673	4.99996074403802\\
-17.1671671671672	5.07323516970612\\
-17.0670670670671	5.1471480411656\\
-16.966966966967	5.22169606562543\\
-16.8668668668669	5.29687574013018\\
-16.7667667667668	5.37268334970997\\
-16.6666666666667	5.44911496559392\\
-16.5665665665666	5.52616644348902\\
-16.4664664664665	5.60383342192624\\
-16.3663663663664	5.68211132067572\\
-16.2662662662663	5.76099533923274\\
-16.1661661661662	5.84048045537627\\
-16.0660660660661	5.92056142380189\\
-15.965965965966	6.00123277483076\\
-15.8658658658659	6.08248881319642\\
-15.7657657657658	6.16432361691098\\
-15.6656656656657	6.24673103621249\\
-15.5655655655656	6.32970469259508\\
-15.4654654654655	6.41323797792344\\
-15.3653653653654	6.4973240536333\\
-15.2652652652653	6.58195585001936\\
-15.1651651651652	6.66712606561228\\
-15.0650650650651	6.75282716664621\\
-14.964964964965	6.83905138661818\\
-14.8648648648649	6.92579072594099\\
-14.7647647647648	7.01303695169072\\
-14.6646646646647	7.10078159745047\\
-14.5645645645646	7.18901596325137\\
-14.4644644644645	7.2777311156123\\
-14.3643643643644	7.36691788767946\\
-14.2642642642643	7.45656687946697\\
-14.1641641641642	7.54666845819959\\
-14.0640640640641	7.63721275875873\\
-13.963963963964	7.72818968423268\\
-13.8638638638639	7.81958890657208\\
-13.7637637637638	7.91139986735156\\
-13.6636636636637	8.00361177863843\\
-13.5635635635636	8.09621362396922\\
-13.4634634634635	8.18919415943487\\
-13.3633633633634	8.28254191487523\\
-13.2632632632633	8.3762451951836\\
-13.1631631631632	8.47029208172184\\
-13.0630630630631	8.56467043384665\\
-12.962962962963	8.6593678905474\\
-12.8628628628629	8.75437187219605\\
-12.7627627627628	8.84966958240942\\
-12.6626626626627	8.94524801002418\\
-12.5625625625626	9.0410939311847\\
-12.4624624624625	9.137193911544\\
-12.3623623623624	9.23353430857785\\
-12.2622622622623	9.33010127401211\\
-12.1621621621622	9.42688075636315\\
-12.0620620620621	9.52385850359148\\
-11.961961961962	9.62102006586807\\
-11.8618618618619	9.71835079845355\\
-11.7617617617618	9.81583586468954\\
-11.6616616616617	9.9134602391021\\
-11.5615615615616	10.0112087106166\\
-11.4614614614615	10.1090658858836\\
-11.3613613613614	10.2070161927152\\
-11.2612612612613	10.3050438836309\\
-11.1611611611612	10.403133039513\\
-11.0610610610611	10.5012675733692\\
-10.960960960961	10.5994312342039\\
-10.8608608608609	10.6976076109942\\
-10.7607607607608	10.795780136773\\
-10.6606606606607	10.8939320928149\\
-10.5605605605606	10.9920466129257\\
-10.4604604604605	11.0901066878342\\
-10.3603603603604	11.1880951696837\\
-10.2602602602603	11.285994776623\\
-10.1601601601602	11.3837880974954\\
-10.0600600600601	11.4814575966232\\
-9.95995995995996	11.578985618687\\
-9.85985985985986	11.6763543936986\\
-9.75975975975976	11.7735460420636\\
-9.65965965965966	11.8705425797347\\
-9.55955955955956	11.9673259234516\\
-9.45945945945946	12.0638778960669\\
-9.35935935935936	12.1601802319555\\
-9.25925925925926	12.2562145825048\\
-9.15915915915916	12.3519625216853\\
-9.05905905905906	12.4474055516973\\
-8.95895895895896	12.5425251086928\\
-8.85885885885886	12.6373025685701\\
-8.75875875875876	12.7317192528381\\
-8.65865865865866	12.8257564345487\\
-8.55855855855856	12.919395344294\\
-8.45845845845846	13.012617176266\\
-8.35835835835836	13.1054030943766\\
-8.25825825825826	13.1977342384348\\
-8.15815815815816	13.2895917303779\\
-8.05805805805806	13.3809566805554\\
-7.95795795795796	13.471810194061\\
-7.85785785785786	13.5621333771109\\
-7.75775775775776	13.6519073434654\\
-7.65765765765766	13.7411132208905\\
-7.55755755755756	13.829732157656\\
-7.45745745745746	13.9177453290691\\
-7.35735735735736	14.0051339440373\\
-7.25725725725726	14.0918792516602\\
-7.15715715715716	14.1779625478451\\
-7.05705705705706	14.2633651819452\\
-6.95695695695696	14.3480685634146\\
-6.85685685685686	14.4320541684788\\
-6.75675675675676	14.5153035468172\\
-6.65665665665666	14.5977983282527\\
-6.55655655655656	14.6795202294476\\
-6.45645645645646	14.7604510605999\\
-6.35635635635636	14.840572732138\\
-6.25625625625626	14.9198672614105\\
-6.15615615615616	14.9983167793666\\
-6.05605605605606	15.075903537225\\
-5.95595595595596	15.1526099131265\\
-5.85585585585586	15.2284184187678\\
-5.75575575575576	15.3033117060128\\
-5.65565565565566	15.3772725734778\\
-5.55555555555556	15.4502839730868\\
-5.45545545545546	15.5223290165942\\
-5.35535535535536	15.5933909820709\\
-5.25525525525526	15.6634533203501\\
-5.15515515515516	15.7324996614302\\
-5.05505505505506	15.8005138208301\\
-4.95495495495495	15.8674798058952\\
-4.85485485485485	15.9333818220482\\
-4.75475475475475	15.9982042789841\\
-4.65465465465465	16.0619317968039\\
-4.55455455455455	16.124549212084\\
-4.45445445445445	16.1860415838795\\
-4.35435435435435	16.2463941996557\\
-4.25425425425425	16.3055925811464\\
-4.15415415415415	16.3636224901345\\
-4.05405405405405	16.4204699341527\\
-3.95395395395395	16.4761211721006\\
-3.85385385385385	16.5305627197742\\
-3.75375375375375	16.5837813553073\\
-3.65365365365365	16.6357641245177\\
-3.55355355355355	16.6864983461596\\
-3.45345345345345	16.7359716170758\\
-3.35335335335335	16.7841718172482\\
-3.25325325325325	16.8310871147441\\
-3.15315315315315	16.8767059705543\\
-3.05305305305305	16.9210171433216\\
-2.95295295295295	16.9640096939552\\
-2.85285285285285	17.0056729901308\\
-2.75275275275275	17.0459967106711\\
-2.65265265265265	17.0849708498065\\
-2.55255255255255	17.1225857213117\\
-2.45245245245245	17.1588319625177\\
-2.35235235235235	17.1937005381945\\
-2.25225225225225	17.2271827443047\\
-2.15215215215215	17.2592702116241\\
-2.05205205205205	17.2899549092279\\
-1.95195195195195	17.3192291478402\\
-1.85185185185185	17.3470855830457\\
-1.75175175175175	17.3735172183596\\
-1.65165165165165	17.3985174081568\\
-1.55155155155155	17.4220798604562\\
-1.45145145145145	17.4441986395601\\
-1.35135135135135	17.4648681685459\\
-1.25125125125125	17.4840832316096\\
-1.15115115115115	17.5018389762598\\
-1.05105105105105	17.5181309153592\\
-0.950950950950951	17.5329549290151\\
-0.850850850850851	17.5463072663156\\
-0.750750750750751	17.5581845469112\\
-0.650650650650651	17.5685837624413\\
-0.550550550550551	17.5775022778038\\
-0.45045045045045	17.5849378322684\\
-0.35035035035035	17.5908885404312\\
-0.25025025025025	17.5953528930118\\
-0.15015015015015	17.5983297574912\\
-0.0500500500500501	17.5998183785906\\
0.0500500500500501	17.5998183785906\\
0.15015015015015	17.5983297574912\\
0.25025025025025	17.5953528930118\\
0.35035035035035	17.5908885404312\\
0.45045045045045	17.5849378322684\\
0.550550550550551	17.5775022778038\\
0.650650650650651	17.5685837624413\\
0.750750750750751	17.5581845469112\\
0.850850850850851	17.5463072663156\\
0.950950950950951	17.5329549290151\\
1.05105105105105	17.5181309153592\\
1.15115115115115	17.5018389762598\\
1.25125125125125	17.4840832316096\\
1.35135135135135	17.4648681685459\\
1.45145145145145	17.4441986395601\\
1.55155155155155	17.4220798604562\\
1.65165165165165	17.3985174081568\\
1.75175175175175	17.3735172183596\\
1.85185185185185	17.3470855830457\\
1.95195195195195	17.3192291478402\\
2.05205205205205	17.2899549092279\\
2.15215215215215	17.2592702116241\\
2.25225225225225	17.2271827443047\\
2.35235235235235	17.1937005381945\\
2.45245245245245	17.1588319625177\\
2.55255255255255	17.1225857213117\\
2.65265265265265	17.0849708498065\\
2.75275275275275	17.0459967106711\\
2.85285285285285	17.0056729901308\\
2.95295295295295	16.9640096939552\\
3.05305305305305	16.9210171433216\\
3.15315315315315	16.8767059705543\\
3.25325325325325	16.8310871147441\\
3.35335335335335	16.7841718172482\\
3.45345345345345	16.7359716170758\\
3.55355355355355	16.6864983461596\\
3.65365365365365	16.6357641245177\\
3.75375375375375	16.5837813553073\\
3.85385385385385	16.5305627197742\\
3.95395395395395	16.4761211721006\\
4.05405405405405	16.4204699341527\\
4.15415415415415	16.3636224901345\\
4.25425425425425	16.3055925811464\\
4.35435435435435	16.2463941996557\\
4.45445445445445	16.1860415838795\\
4.55455455455455	16.124549212084\\
4.65465465465465	16.0619317968039\\
4.75475475475475	15.9982042789841\\
4.85485485485485	15.9333818220482\\
4.95495495495495	15.8674798058952\\
5.05505505505506	15.8005138208301\\
5.15515515515516	15.7324996614302\\
5.25525525525526	15.6634533203501\\
5.35535535535536	15.5933909820709\\
5.45545545545546	15.5223290165942\\
5.55555555555556	15.4502839730868\\
5.65565565565566	15.3772725734778\\
5.75575575575576	15.3033117060128\\
5.85585585585586	15.2284184187678\\
5.95595595595596	15.1526099131265\\
6.05605605605606	15.075903537225\\
6.15615615615616	14.9983167793666\\
6.25625625625626	14.9198672614105\\
6.35635635635636	14.840572732138\\
6.45645645645646	14.7604510605999\\
6.55655655655656	14.6795202294476\\
6.65665665665666	14.5977983282527\\
6.75675675675676	14.5153035468172\\
6.85685685685686	14.4320541684788\\
6.95695695695696	14.3480685634146\\
7.05705705705706	14.2633651819452\\
7.15715715715716	14.1779625478451\\
7.25725725725726	14.0918792516602\\
7.35735735735736	14.0051339440373\\
7.45745745745746	13.9177453290691\\
7.55755755755756	13.829732157656\\
7.65765765765766	13.7411132208905\\
7.75775775775776	13.6519073434654\\
7.85785785785786	13.5621333771109\\
7.95795795795796	13.471810194061\\
8.05805805805806	13.3809566805554\\
8.15815815815816	13.2895917303779\\
8.25825825825826	13.1977342384348\\
8.35835835835836	13.1054030943766\\
8.45845845845846	13.012617176266\\
8.55855855855856	12.919395344294\\
8.65865865865866	12.8257564345487\\
8.75875875875876	12.7317192528381\\
8.85885885885886	12.6373025685701\\
8.95895895895896	12.5425251086928\\
9.05905905905906	12.4474055516973\\
9.15915915915916	12.3519625216853\\
9.25925925925926	12.2562145825048\\
9.35935935935936	12.1601802319555\\
9.45945945945946	12.0638778960669\\
9.55955955955956	11.9673259234516\\
9.65965965965966	11.8705425797347\\
9.75975975975976	11.7735460420636\\
9.85985985985986	11.6763543936986\\
9.95995995995996	11.578985618687\\
10.0600600600601	11.4814575966232\\
10.1601601601602	11.3837880974954\\
10.2602602602603	11.285994776623\\
10.3603603603604	11.1880951696837\\
10.4604604604605	11.0901066878342\\
10.5605605605606	10.9920466129257\\
10.6606606606607	10.8939320928149\\
10.7607607607608	10.795780136773\\
10.8608608608609	10.6976076109942\\
10.960960960961	10.5994312342039\\
11.0610610610611	10.5012675733692\\
11.1611611611612	10.403133039513\\
11.2612612612613	10.3050438836309\\
11.3613613613614	10.2070161927152\\
11.4614614614615	10.1090658858836\\
11.5615615615616	10.0112087106166\\
11.6616616616617	9.9134602391021\\
11.7617617617618	9.81583586468954\\
11.8618618618619	9.71835079845355\\
11.961961961962	9.62102006586807\\
12.0620620620621	9.52385850359148\\
12.1621621621622	9.42688075636315\\
12.2622622622623	9.33010127401211\\
12.3623623623624	9.23353430857785\\
12.4624624624625	9.137193911544\\
12.5625625625626	9.0410939311847\\
12.6626626626627	8.94524801002418\\
12.7627627627628	8.84966958240942\\
12.8628628628629	8.75437187219605\\
12.962962962963	8.6593678905474\\
13.0630630630631	8.56467043384665\\
13.1631631631632	8.47029208172184\\
13.2632632632633	8.3762451951836\\
13.3633633633634	8.28254191487523\\
13.4634634634635	8.18919415943487\\
13.5635635635636	8.09621362396922\\
13.6636636636637	8.00361177863843\\
13.7637637637638	7.91139986735156\\
13.8638638638639	7.81958890657208\\
13.963963963964	7.72818968423268\\
14.0640640640641	7.63721275875873\\
14.1641641641642	7.54666845819959\\
14.2642642642643	7.45656687946697\\
14.3643643643644	7.36691788767946\\
14.4644644644645	7.2777311156123\\
14.5645645645646	7.18901596325137\\
14.6646646646647	7.10078159745047\\
14.7647647647648	7.01303695169072\\
14.8648648648649	6.92579072594099\\
14.964964964965	6.83905138661818\\
15.0650650650651	6.75282716664621\\
15.1651651651652	6.66712606561228\\
15.2652652652653	6.58195585001936\\
15.3653653653654	6.4973240536333\\
15.4654654654655	6.41323797792344\\
15.5655655655656	6.32970469259508\\
15.6656656656657	6.24673103621249\\
15.7657657657658	6.16432361691098\\
15.8658658658659	6.08248881319642\\
15.965965965966	6.00123277483076\\
16.0660660660661	5.92056142380189\\
16.1661661661662	5.84048045537627\\
16.2662662662663	5.76099533923274\\
16.3663663663664	5.68211132067572\\
16.4664664664665	5.60383342192624\\
16.5665665665666	5.52616644348902\\
16.6666666666667	5.44911496559392\\
16.7667667667668	5.37268334970997\\
16.8668668668669	5.29687574013018\\
16.966966966967	5.22169606562543\\
17.0670670670671	5.1471480411656\\
17.1671671671672	5.07323516970612\\
17.2672672672673	4.99996074403802\\
17.3673673673674	4.92732784869987\\
17.4674674674675	4.85533936194955\\
17.5675675675676	4.78399795779399\\
17.6676676676677	4.71330610807525\\
17.7677677677678	4.64326608461076\\
17.8678678678679	4.57387996138604\\
17.967967967968	4.50514961679799\\
18.0680680680681	4.43707673594682\\
18.1681681681682	4.36966281297475\\
18.2682682682683	4.30290915344969\\
18.3683683683684	4.23681687679183\\
18.4684684684685	4.17138691874151\\
18.5685685685686	4.10662003386626\\
18.6686686686687	4.04251679810529\\
18.7687687687688	3.97907761134962\\
18.8688688688689	3.91630270005573\\
18.968968968969	3.85419211989127\\
19.0690690690691	3.79274575841064\\
19.1691691691692	3.7319633377589\\
19.2692692692693	3.67184441740198\\
19.3693693693694	3.61238839688158\\
19.4694694694695	3.55359451859281\\
19.5695695695696	3.49546187058299\\
19.6696696696697	3.43798938936967\\
19.7697697697698	3.38117586277622\\
19.8698698698699	3.32501993278338\\
19.96996996997	3.26952009839489\\
20.0700700700701	3.21467471851561\\
20.1701701701702	3.16048201484052\\
20.2702702702703	3.10694007475287\\
20.3703703703704	3.05404685422997\\
20.4704704704705	3.00180018075488\\
20.5705705705706	2.95019775623268\\
20.6706706706707	2.89923715990948\\
20.7707707707708	2.84891585129284\\
20.8708708708709	2.79923117307216\\
20.970970970971	2.75018035403731\\
21.0710710710711	2.70176051199441\\
21.1711711711712	2.65396865667699\\
21.2712712712713	2.60680169265145\\
21.3713713713714	2.56025642221529\\
21.4714714714715	2.51432954828676\\
21.5715715715716	2.4690176772848\\
21.6716716716717	2.42431732199776\\
21.7717717717718	2.38022490443988\\
21.8718718718719	2.33673675869416\\
21.971971971972	2.2938491337405\\
22.0720720720721	2.25155819626799\\
22.1721721721722	2.2098600334701\\
22.2722722722723	2.16875065582185\\
22.3723723723724	2.12822599983775\\
22.4724724724725	2.08828193080957\\
22.5725725725726	2.04891424552287\\
22.6726726726727	2.0101186749514\\
22.7727727727728	1.97189088692837\\
22.8728728728729	1.93422648879368\\
22.972972972973	1.89712103001628\\
23.0730730730731	1.8605700047908\\
23.1731731731732	1.82456885460762\\
23.2732732732733	1.78911297079558\\
23.3733733733734	1.75419769703661\\
23.4734734734735	1.71981833185162\\
23.5735735735736	1.68597013105679\\
23.6736736736737	1.65264831018979\\
23.7737737737738	1.61984804690515\\
23.8738738738739	1.58756448333834\\
23.973973973974	1.55579272843779\\
24.0740740740741	1.52452786026452\\
24.1741741741742	1.49376492825872\\
24.2742742742743	1.46349895547287\\
24.3743743743744	1.43372494077098\\
24.4744744744745	1.40443786099339\\
24.5745745745746	1.375632673087\\
24.6746746746747	1.34730431620018\\
24.7747747747748	1.3194477137425\\
24.8748748748749	1.29205777540849\\
24.974974974975	1.26512939916557\\
25.0750750750751	1.23865747320553\\
25.1751751751752	1.2126368778597\\
25.2752752752753	1.18706248747723\\
25.3753753753754	1.16192917226666\\
25.4754754754755	1.13723180010031\\
25.5755755755756	1.11296523828165\\
25.6756756756757	1.08912435527534\\
25.7757757757758	1.06570402239998\\
25.8758758758759	1.04269911548341\\
25.975975975976	1.0201045164807\\
26.0760760760761	0.99791511505465\\
26.1761761761762	0.976125810118873\\
26.2762762762763	0.954731511343568\\
26.3763763763764	0.933727140623928\\
26.4764764764765	0.913107633511332\\
26.5765765765766	0.892867940607398\\
26.6766766766767	0.873003028921016\\
26.7767767767768	0.853507883188493\\
26.8768768768769	0.834377507156975\\
26.976976976977	0.815606924831305\\
27.0770770770771	0.79719118168453\\
27.1771771771772	0.779125345832233\\
27.2772772772773	0.761404509170944\\
27.3773773773774	0.744023788480854\\
27.4774774774775	0.726978326493092\\
27.5775775775776	0.710263292921835\\
27.6776776776777	0.693873885461536\\
27.7777777777778	0.677805330749564\\
27.8778778778779	0.66205288529459\\
27.977977977978	0.646611836371\\
28.0780780780781	0.631477502879726\\
28.1781781781782	0.616645236175803\\
28.2782782782783	0.602110420863032\\
28.3783783783784	0.587868475556115\\
28.4784784784785	0.573914853610646\\
28.5785785785786	0.560245043821347\\
28.6786786786787	0.546854571088956\\
28.7787787787788	0.533738997056165\\
28.8788788788789	0.520893920713036\\
28.978978978979	0.508314978972322\\
29.0790790790791	0.495997847215118\\
29.1791791791792	0.48393823980729\\
29.2792792792793	0.472131910587122\\
29.3793793793794	0.460574653324644\\
29.4794794794795	0.449262302153084\\
29.5795795795796	0.438190731972929\\
29.6796796796797	0.427355858829039\\
29.7797797797798	0.416753640261303\\
29.8798798798799	0.406380075629314\\
29.97997997998	0.396231206411528\\
30.0800800800801	0.386303116479408\\
30.1801801801802	0.376591932347021\\
30.2802802802803	0.367093823396597\\
30.3803803803804	0.357805002080515\\
30.4804804804805	0.348721724100225\\
30.5805805805806	0.339840288562594\\
30.6806806806807	0.331157038114162\\
30.7807807807808	0.322668359053812\\
30.8808808808809	0.31437068142434\\
30.980980980981	0.306260479083419\\
31.0810810810811	0.298334269754454\\
31.1811811811812	0.290588615057812\\
31.2812812812813	0.283020120522929\\
31.3813813813814	0.275625435581768\\
31.4814814814815	0.268401253544132\\
31.5815815815816	0.261344311555301\\
31.6816816816817	0.254451390536484\\
31.7817817817818	0.247719315108576\\
31.8818818818819	0.241144953499672\\
31.981981981982	0.234725217436837\\
32.0820820820821	0.22845706202259\\
32.1821821821822	0.222337485596574\\
32.2822822822823	0.216363529582873\\
32.3823823823824	0.21053227832344\\
32.4824824824825	0.204840858898081\\
32.5825825825826	0.199286440931466\\
32.6826826826827	0.193866236387588\\
32.7827827827828	0.188577499352136\\
32.8828828828829	0.183417525803198\\
32.982982982983	0.178383653370747\\
33.0830830830831	0.173473261085314\\
33.1831831831832	0.168683769116292\\
33.2832832832833	0.164012638500266\\
33.3833833833834	0.159457370859795\\
33.4834834834835	0.155015508113041\\
33.5835835835836	0.150684632174654\\
33.6836836836837	0.146462364648285\\
33.7837837837838	0.142346366511146\\
33.8838838838839	0.138334337790961\\
33.983983983984	0.134424017235712\\
34.0840840840841	0.130613181976535\\
34.1841841841842	0.126899647184121\\
34.2842842842843	0.123281265718999\\
34.3843843843844	0.119755927776021\\
34.4844844844845	0.116321560523418\\
34.5845845845846	0.112976127736744\\
34.6846846846847	0.109717629428052\\
34.7847847847848	0.106544101470602\\
34.8848848848849	0.103453615219449\\
34.984984984985	0.100444277128188\\
35.0850850850851	0.0975142283621819\\
35.1851851851852	0.0946616444085526\\
35.2852852852853	0.0918847346832355\\
35.3853853853854	0.089181742135371\\
35.4854854854855	0.0865509428493145\\
35.5855855855856	0.0839906456445301\\
35.6856856856857	0.0814991916736334\\
35.7857857857858	0.0790749540188357\\
35.8858858858859	0.0767163372870418\\
35.985985985986	0.0744217772038405\\
36.0860860860861	0.0721897402066271\\
36.1861861861862	0.0700187230370848\\
36.2862862862863	0.0679072523332475\\
36.3863863863864	0.0658538842213639\\
36.4864864864865	0.0638572039077676\\
36.5865865865866	0.0619158252709626\\
36.6866866866867	0.0600283904541162\\
36.7867867867868	0.0581935694581557\\
36.8868868868869	0.0564100597356489\\
36.986986986987	0.0546765857856524\\
37.0870870870871	0.0529918987496952\\
37.1871871871872	0.0513547760090706\\
37.2872872872873	0.0497640207835921\\
37.3873873873874	0.0482184617319718\\
37.4874874874875	0.0467169525539695\\
37.5875875875876	0.0452583715944576\\
37.6876876876877	0.0438416214495384\\
37.7877877877878	0.0424656285748482\\
37.8878878878879	0.0411293428961744\\
37.987987987988	0.0398317374225089\\
38.0880880880881	0.0385718078616531\\
38.1881881881882	0.0373485722384879\\
38.2882882882883	0.0361610705160136\\
38.3883883883884	0.0350083642192628\\
38.4884884884885	0.0338895360621817\\
38.5885885885886	0.0328036895775726\\
38.6886886886887	0.0317499487501851\\
38.7887887887888	0.0307274576530372\\
38.8888888888889	0.0297353800870457\\
38.988988988989	0.0287728992240389\\
39.0890890890891	0.0278392172532215\\
39.1891891891892	0.0269335550311556\\
39.2892892892893	0.0260551517353202\\
39.3893893893894	0.0252032645213043\\
39.4894894894895	0.0243771681836888\\
39.5895895895896	0.0235761548206637\\
39.6896896896897	0.0227995335024272\\
39.7897897897898	0.0220466299434088\\
39.8898898898899	0.0213167861783529\\
39.98998998999	0.020609360242298\\
40.0900900900901	0.0199237258544833\\
40.1901901901902	0.0192592721062093\\
40.2902902902903	0.0186154031526777\\
40.3903903903904	0.0179915379088314\\
40.4904904904905	0.0173871097492134\\
40.5905905905906	0.0168015662118602\\
40.6906906906907	0.0162343687062415\\
40.7907907907908	0.0156849922252566\\
40.8908908908909	0.0151529250612943\\
40.990990990991	0.0146376685263606\\
41.0910910910911	0.0141387366762767\\
41.1911911911912	0.0136556560389466\\
41.2912912912913	0.0131879653466907\\
41.3913913913914	0.0127352152726418\\
41.4914914914915	0.0122969681711943\\
41.5915915915916	0.0118727978224996\\
41.6916916916917	0.0114622891809937\\
41.7917917917918	0.0110650381279454\\
41.8918918918919	0.0106806512280092\\
41.991991991992	0.010308745489766\\
42.0920920920921	0.00994894813023214\\
42.1921921921922	0.00960089634331778\\
42.2922922922923	0.00926423707221086\\
42.3923923923924	0.0089386267856644\\
42.4924924924925	0.00862373125816143\\
42.5925925925926	0.00831922535393144\\
42.6926926926927	0.00802479281479058\\
42.7927927927928	0.00774012605177678\\
42.8928928928929	0.00746492594054935\\
42.992992992993	0.00719890162052228\\
43.0930930930931	0.00694177029769847\\
43.1931931931932	0.00669325705117179\\
43.2932932932933	0.00645309464326277\\
43.3933933933934	0.00622102333325249\\
43.4934934934935	0.00599679069467906\\
43.5935935935936	0.00578015143615947\\
43.6936936936937	0.00557086722569968\\
43.7937937937938	0.00536870651845469\\
43.8938938938939	0.00517344438789945\\
43.993993993994	0.00498486236037177\\
44.0940940940941	0.00480274825294689\\
44.1941941941942	0.00462689601460352\\
44.2942942942943	0.0044571055706403\\
44.3943943943944	0.00429318267030165\\
44.4944944944945	0.00413493873757133\\
44.5945945945946	0.00398219072509176\\
44.6946946946947	0.00383476097116684\\
44.7947947947948	0.00369247705980604\\
44.8948948948949	0.0035551716837669\\
44.994994994995	0.00342268251055332\\
45.0950950950951	0.00329485205132659\\
45.1951951951952	0.00317152753268615\\
45.2952952952953	0.00305256077127702\\
45.3953953953954	0.00293780805118084\\
45.4954954954955	0.00282713000404714\\
45.5955955955956	0.00272039149192201\\
45.6956956956957	0.002617461492731\\
45.7957957957958	0.0025182129883732\\
45.8958958958959	0.00242252285538374\\
45.995995995996	0.00233027175812187\\
46.0960960960961	0.00224134404444221\\
46.1961961961962	0.00215562764380639\\
46.2962962962963	0.00207301396779331\\
46.3963963963964	0.00199339781296561\\
46.4964964964965	0.00191667726605093\\
46.5965965965966	0.00184275361139605\\
46.6966966966967	0.00177153124065318\\
46.7967967967968	0.00170291756465684\\
46.8968968968969	0.00163682292745133\\
46.996996996997	0.00157316052242787\\
47.0970970970971	0.00151184631053185\\
47.1971971971972	0.00145279894050035\\
47.2972972972973	0.00139593967109069\\
47.3973973973974	0.00134119229526108\\
47.4974974974975	0.00128848306626492\\
47.5975975975976	0.00123774062562056\\
47.6976976976977	0.00118889593291883\\
47.7977977977978	0.00114188219743101\\
47.8978978978979	0.00109663481148037\\
47.997997997998	0.00105309128554084\\
48.0980980980981	0.0010111911850268\\
48.1981981981982	0.000970876068738363\\
48.2982982982983	0.000932089428927083\\
48.3983983983984	0.000894776632947465\\
48.4984984984985	0.00085888486645993\\
48.5985985985986	0.000824363078151659\\
48.6986986986987	0.000791161925941933\\
48.7987987987988	0.000759233724639237\\
48.8988988988989	0.000728532395017747\\
48.998998998999	0.000699013414281412\\
49.0990990990991	0.000670633767884217\\
49.1991991991992	0.000643351902675775\\
49.2992992992993	0.000617127681341812\\
49.3993993993994	0.000591922338109637\\
49.4994994994995	0.000567698435689168\\
49.5995995995996	0.000544419823420536\\
49.6996996996997	0.000522051596599788\\
49.7997997997998	0.000500560056954716\\
49.8998998998999	0.000479912674243277\\
50	0.000460078048947578\\
};
\addlegendentry{Distribution}

\addplot [color=red, only marks, mark size=2.5pt, mark=*, mark options={solid, fill=red, red}]
  table[row sep=crcr]{%
0	17.6000044650816\\
35	0.0999997600913489\\
};
\addlegendentry{Data}

\end{axis}
\end{tikzpicture}%
\end{figure}
This would give us the corrected values of intensities.
However, doing so does not affect the readings near the minimum (\(0\) in the plot).
This does not explain why the zero is not present in any of the readings, and how say that the minimum value is the brewster angle.

Light that is reflected from a surface is highly polarized when the angle between the reflected ray and refracted ray is close to \(90^\circ\).
By adding a polarizing filter to the incoming light, the amount of reflected light can be controlled by changed by rotating the polarizing filter.
If the filter is not properly alligned, a small component of the light polarized in the direction the reflected beam gets through which means there is a non zero reading at brewster's angle.
We see the same in the figure below.
\begin{figure}[H]
    \centering
    \setlength\fwidth{0.75\textwidth}
    \setlength\fheight{0.5\textwidth}
    % This file was created by matlab2tikz.
%
\definecolor{mycolor1}{rgb}{0.00000,0.44700,0.74100}%
\definecolor{mycolor2}{rgb}{0.85000,0.32500,0.09800}%
\definecolor{mycolor3}{rgb}{0.12941,0.12941,0.12941}%
%
\begin{tikzpicture}

\begin{axis}[%
width=0.951\fwidth,
height=\fheight,
at={(0\fwidth,0\fheight)},
scale only axis,
xmin=0,
xmax=500,
xlabel style={font=\color{mycolor3}},
xlabel={Polarizer angle (in degrees)},
ymin=0,
ymax=1500,
ylabel style={font=\color{mycolor3}},
ylabel={Intensity (in $\mu$A)},
axis background/.style={fill=white},
xmajorgrids,
ymajorgrids,
legend style={legend cell align=left, align=left}
]
\addplot [color=mycolor1]
  table[row sep=crcr]{%
-50	297.752952021078\\
-49.4494494494494	308.855134411315\\
-48.8988988988989	320.102464743158\\
-48.3483483483483	331.490789465135\\
-47.7977977977978	343.015902957646\\
-47.2472472472472	354.673549086062\\
-46.6966966966967	366.45942277248\\
-46.1461461461461	378.369171585572\\
-45.5955955955956	390.398397347879\\
-45.045045045045	402.542657760038\\
-44.4944944944945	414.797468041286\\
-43.9439439439439	427.158302585656\\
-43.3933933933934	439.62059663324\\
-42.8428428428428	452.179747955942\\
-42.2922922922923	464.831118557013\\
-41.7417417417417	477.570036383849\\
-41.1911911911912	490.391797053336\\
-40.6406406406406	503.291665589145\\
-40.0900900900901	516.26487817032\\
-39.5395395395395	529.306643890515\\
-38.988988988989	542.41214652726\\
-38.4384384384384	555.576546320521\\
-37.8878878878879	568.79498176001\\
-37.3373373373373	582.062571380496\\
-36.7867867867868	595.37441556449\\
-36.2362362362362	608.725598351633\\
-35.6856856856857	622.111189254143\\
-35.1351351351351	635.526245077573\\
-34.5845845845846	648.965811746313\\
-34.034034034034	662.424926133093\\
-33.4834834834835	675.89861789182\\
-32.9329329329329	689.381911293093\\
-32.3823823823824	702.869827061691\\
-31.8318318318318	716.357384215408\\
-31.2812812812813	729.839601904453\\
-30.7307307307307	743.311501250866\\
-30.1801801801802	756.768107187168\\
-29.6296296296296	770.204450293617\\
-29.0790790790791	783.615568633374\\
-28.5285285285285	796.996509584923\\
-27.977977977978	810.342331671007\\
-27.4274274274274	823.648106383497\\
-26.8768768768769	836.90892000345\\
-26.3263263263263	850.119875415701\\
-25.7757757757758	863.276093917329\\
-25.2252252252252	876.372717019342\\
-24.6746746746747	889.404908240852\\
-24.1241241241241	902.367854895165\\
-23.5735735735736	915.256769867067\\
-23.023023023023	928.066893380667\\
-22.4724724724725	940.793494757148\\
-21.9219219219219	953.431874161761\\
-21.3713713713714	965.977364339459\\
-20.8208208208208	978.425332338446\\
-20.2702702702703	990.771181221115\\
-19.7197197197197	1003.01035176165\\
-19.1691691691692	1015.13832412972\\
-18.6186186186186	1027.15061955961\\
-18.0680680680681	1039.0428020042\\
-17.5175175175175	1050.81047977315\\
-16.966966966967	1062.44930715474\\
-16.4164164164164	1073.95498602069\\
-15.8658658658659	1085.32326741342\\
-15.3153153153153	1096.54995311518\\
-14.7647647647648	1107.6308971984\\
-14.2142142142142	1118.56200755675\\
-13.6636636636637	1129.33924741634\\
-13.1131131131131	1139.95863682644\\
-12.5625625625626	1150.41625412927\\
-12.012012012012	1160.70823740823\\
-11.4614614614615	1170.83078591407\\
-10.9109109109109	1180.78016146849\\
-10.3603603603604	1190.55268984461\\
-9.80980980980981	1200.14476212382\\
-9.25925925925926	1209.55283602857\\
-8.70870870870871	1218.77343723044\\
-8.15815815815816	1227.80316063324\\
-7.6076076076076	1236.63867163047\\
-7.05705705705706	1245.27670733676\\
-6.5065065065065	1253.71407779279\\
-5.95595595595596	1261.9476671434\\
-5.4054054054054	1269.97443478819\\
-4.85485485485486	1277.79141650438\\
-4.3043043043043	1285.39572554153\\
-3.75375375375376	1292.78455368755\\
-3.2032032032032	1299.95517230577\\
-2.65265265265266	1306.9049333426\\
-2.1021021021021	1313.63127030542\\
-1.55155155155155	1320.1316992104\\
-1.001001001001	1326.40381949977\\
-0.450450450450454	1332.44531492839\\
0.1001001001001	1338.25395441905\\
0.650650650650654	1343.82759288644\\
1.2012012012012	1349.16417202929\\
1.75175175175175	1354.26172109051\\
2.3023023023023	1359.11835758493\\
2.85285285285286	1363.73228799452\\
3.4034034034034	1368.10180843075\\
3.95395395395396	1372.22530526375\\
4.5045045045045	1376.10125571825\\
5.05505505505506	1379.72822843596\\
5.6056056056056	1383.10488400411\\
6.15615615615616	1386.22997545009\\
6.7067067067067	1389.10234870198\\
7.25725725725726	1391.72094301472\\
7.80780780780781	1394.08479136183\\
8.35835835835836	1396.19302079253\\
8.90890890890891	1398.04485275411\\
9.45945945945946	1399.63960337945\\
10.01001001001	1400.97668373958\\
10.5605605605606	1402.05560006113\\
11.1111111111111	1402.87595390871\\
11.6616616616617	1403.43744233204\\
12.2122122122122	1403.73985797781\\
12.7627627627628	1403.78308916628\\
13.3133133133133	1403.56711993252\\
13.8638638638639	1403.09203003227\\
14.4144144144144	1402.35799491254\\
14.964964964965	1401.36528564678\\
15.5155155155155	1400.11426883478\\
16.0660660660661	1398.60540646731\\
16.6166166166166	1396.83925575549\\
17.1671671671672	1394.81646892501\\
17.7177177177177	1392.5377929753\\
18.2682682682683	1390.00406940363\\
18.8188188188188	1387.21623389437\\
19.3693693693694	1384.17531597347\\
19.9199199199199	1380.88243862821\\
20.4704704704705	1377.33881789256\\
21.021021021021	1373.54576239804\\
21.5715715715716	1369.5046728905\\
22.1221221221221	1365.21704171282\\
22.6726726726727	1360.68445225377\\
23.2232232232232	1355.90857836335\\
23.7737737737738	1350.89118373457\\
24.3243243243243	1345.63412125218\\
24.8748748748749	1340.13933230841\\
25.4254254254254	1334.408846086\\
25.975975975976	1328.44477880887\\
26.5265265265265	1322.24933296061\\
27.0770770770771	1315.8247964711\\
27.6276276276276	1309.17354187162\\
28.1781781781782	1302.29802541868\\
28.7287287287287	1295.20078618694\\
29.2792792792793	1287.88444513156\\
29.8298298298298	1280.3517041203\\
30.3803803803804	1272.60534493573\\
30.9309309309309	1264.64822824794\\
31.4814814814815	1256.48329255815\\
32.032032032032	1248.11355311349\\
32.5825825825826	1239.54210079354\\
33.1331331331331	1230.77210096887\\
33.6836836836837	1221.80679233208\\
34.2342342342342	1212.64948570183\\
34.7847847847848	1203.3035628001\\
35.3353353353353	1193.77247500342\\
35.8858858858859	1184.05974206827\\
36.4364364364364	1174.16895083123\\
36.986986986987	1164.10375388448\\
37.5375375375375	1153.86786822683\\
38.0880880880881	1143.46507389111\\
38.6386386386386	1132.89921254824\\
39.1891891891892	1122.17418608849\\
39.7397397397397	1111.29395518057\\
40.2902902902903	1100.26253780902\\
40.8408408408408	1089.08400779031\\
41.3913913913914	1077.76249326852\\
41.9419419419419	1066.30217519074\\
42.4924924924925	1054.70728576317\\
43.043043043043	1042.98210688814\\
43.5935935935936	1031.13096858287\\
44.1441441441441	1019.15824738041\\
44.6946946946947	1007.06836471344\\
45.2452452452452	994.865785281423\\
45.7957957957958	982.5550154019\\
46.3463463463464	970.140601346271\\
46.8968968968969	957.627127660939\\
47.4474474474474	945.019215474251\\
47.997997997998	932.321520789958\\
48.5485485485486	919.538732767789\\
49.0990990990991	906.675571991782\\
49.6496496496496	893.736788726995\\
50.2002002002002	880.727161165293\\
50.7507507507508	867.651493660787\\
51.3013013013013	854.514614955622\\
51.8518518518518	841.321376396764\\
52.4024024024024	828.076650144432\\
52.952952952953	814.785327372856\\
53.5035035035035	801.452316463982\\
54.054054054054	788.082541194875\\
54.6046046046046	774.680938919386\\
55.1551551551552	761.252458744834\\
55.7057057057057	747.802059704334\\
56.2562562562563	734.334708925474\\
56.8068068068068	720.855379795971\\
57.3573573573574	707.369050127071\\
57.9079079079079	693.880700315256\\
58.4584584584585	680.395311503034\\
59.009009009009	666.917863739436\\
59.5595595595596	653.453334140929\\
60.1101101101101	640.006695053401\\
60.6606606606607	626.582912215917\\
61.2112112112112	613.186942926893\\
61.7617617617618	599.823734213435\\
62.3123123123123	586.498221004416\\
62.8628628628629	573.215324308055\\
63.4134134134134	559.979949394627\\
63.963963963964	546.796983984971\\
64.5145145145145	533.671296445518\\
65.0650650650651	520.607733990423\\
65.6156156156156	507.611120891532\\
66.1661661661662	494.686256696812\\
66.7167167167167	481.837914457924\\
67.2672672672673	469.070838967563\\
67.8178178178178	456.389745007252\\
68.3683683683684	443.799315606192\\
68.9189189189189	431.304200311881\\
69.4694694694695	418.909013473058\\
70.02002002002	406.618332535658\\
70.5705705705706	394.436696352408\\
71.1211211211211	382.368603506643\\
71.6716716716717	370.418510651041\\
72.2222222222222	358.590830861798\\
72.7727727727728	346.889932008919\\
73.3233233233233	335.320135143196\\
73.8738738738739	323.88571290048\\
74.4244244244244	312.590887923826\\
74.974974974975	301.439831304113\\
75.5255255255255	290.436661039671\\
76.0760760760761	279.585440515565\\
76.6266266266266	268.890177003\\
77.1771771771772	258.354820179471\\
77.7277277277277	247.983260670183\\
78.2782782782783	237.779328611259\\
78.8288288288288	227.746792235318\\
79.3793793793794	217.889356479888\\
79.9299299299299	208.2106616192\\
80.4804804804805	198.714281919867\\
81.031031031031	189.403724320931\\
81.5815815815816	180.282427138782\\
82.1321321321321	171.353758797414\\
82.6826826826827	162.621016584482\\
83.2332332332332	154.087425433658\\
83.7837837837838	145.756136733673\\
84.3343343343343	137.630227164542\\
84.8848848848849	129.712697561368\\
85.4354354354354	122.006471806147\\
85.985985985986	114.51439574802\\
86.5365365365365	107.239236152308\\
87.0870870870871	100.183679678776\\
87.6376376376376	93.3503318894632\\
88.1881881881882	86.7417162864745\\
88.7387387387387	80.3602733800711\\
89.2892892892893	74.2083597874\\
89.8398398398398	68.2882473622282\\
90.3903903903904	62.6021223559545\\
90.9409409409409	57.1520846102471\\
91.4914914914915	51.9401467815878\\
92.0420420420421	46.968233598014\\
92.5925925925926	42.2381811483333\\
93.1431431431431	37.7517362040633\\
93.6936936936937	33.5105555743728\\
94.2442442442442	29.5162054942281\\
94.7947947947948	25.7701610459965\\
95.3453453453454	22.2738056147099\\
95.8958958958959	19.0284303771913\\
96.4464464464465	16.0352338252343\\
96.996996996997	13.2953213230056\\
97.5475475475475	10.8097046988484\\
98.0980980980981	8.57930187161839\\
98.6486486486486	6.60493651170397\\
99.1991991991992	4.88733773685167\\
99.7497497497498	3.42713984290886\\
100.3003003003	2.22488206958128\\
100.850850850851	1.28100840129923\\
101.401401401401	0.595867403255308\\
101.951951951952	0.169712092682632\\
102.502502502502	0.00269984541750449\\
103.053053053053	0.0948923377816774\\
103.603603603604	0.446255523805999\\
104.154154154154	1.05665964780312\\
104.704704704705	1.92587929228529\\
105.255255255255	3.05359346120931\\
105.805805805806	4.43938569851787\\
106.356356356356	6.08274424193283\\
106.906906906907	7.98306221194764\\
107.457457457457	10.1396378359395\\
108.008008008008	12.5516747073309\\
108.558558558559	15.2182820796958\\
109.109109109109	18.1384751957054\\
109.65965965966	21.311175650791\\
110.21021021021	24.7352117913912\\
110.760760760761	28.409319147632\\
111.311311311311	32.332140900292\\
111.861861861862	36.5022283818568\\
112.412412412412	40.9180416115057\\
112.962962962963	45.5779498638137\\
113.513513513514	50.4802322709654\\
114.064064064064	55.6230784582563\\
114.614614614615	61.0045892126576\\
115.165165165165	66.6227771841696\\
115.715715715716	72.475567619742\\
116.266266266266	78.5607991294621\\
116.816816816817	84.8762244847393\\
117.367367367367	91.4195114481866\\
117.917917917918	98.1882436349075\\
118.468468468468	105.179921404834\\
119.019019019019	112.391962785833\\
119.56956956957	119.821704427206\\
120.12012012012	127.466402583246\\
120.670670670671	135.323234126476\\
121.221221221221	143.389297590218\\
121.771771771772	151.661614240072\\
122.322322322322	160.137129173941\\
122.872872872873	168.812712450189\\
123.423423423423	177.685160243496\\
123.973973973974	186.751196028012\\
124.524524524525	196.007471787349\\
125.075075075075	205.450569250975\\
125.625625625626	215.077001156568\\
126.176176176176	224.883212537816\\
126.726726726727	234.865582037248\\
127.277277277277	245.020423243575\\
127.827827827828	255.343986053045\\
128.378378378378	265.832458054344\\
128.928928928929	276.481965936466\\
129.479479479479	287.288576919115\\
130.03003003003	298.24830020504\\
130.580580580581	309.357088453812\\
131.131131131131	320.610839276473\\
131.681681681682	332.005396750513\\
132.232232232232	343.536552954621\\
132.782782782783	355.200049522651\\
133.333333333333	366.991579216178\\
133.883883883884	378.906787515145\\
134.434434434434	390.941274225953\\
134.984984984985	403.090595106406\\
135.535535535536	415.350263506959\\
136.086086086086	427.715752027575\\
136.636636636637	440.182494189674\\
137.187187187187	452.745886122496\\
137.737737737738	465.401288263272\\
138.288288288288	478.144027070583\\
138.838838838839	490.969396750251\\
139.389389389389	503.872660993177\\
139.93993993994	516.849054724393\\
140.49049049049	529.893785862783\\
141.041041041041	543.002037090761\\
141.591591591592	556.168967633267\\
142.142142142142	569.38971504542\\
142.692692692693	582.659397008205\\
143.243243243243	595.973113131443\\
143.793793793794	609.325946763488\\
144.344344344344	622.712966806904\\
144.894894894895	636.129229539477\\
145.445445445445	649.569780439896\\
145.995995995996	663.029656017422\\
146.546546546547	676.503885644854\\
147.097097097097	689.987493394171\\
147.647647647648	703.475499874067\\
148.198198198198	716.962924068824\\
148.748748748749	730.44478517776\\
149.299299299299	743.91610445459\\
149.84984984985	757.371907046061\\
150.4004004004	770.807223829096\\
150.950950950951	784.217093245872\\
151.501501501502	797.596563136081\\
152.052052052052	810.940692565724\\
152.602602602603	824.244553651757\\
153.153153153153	837.503233381917\\
153.703703703704	850.711835429079\\
154.254254254254	863.865481959402\\
154.804804804805	876.959315433693\\
155.355355355355	889.988500401256\\
155.905905905906	902.948225285583\\
156.456456456456	915.833704161231\\
157.007007007007	928.640178521243\\
157.557557557558	941.362919034398\\
158.108108108108	953.99722729174\\
158.658658658659	966.538437541655\\
159.209209209209	978.981918412898\\
159.75975975976	991.323074624921\\
160.31031031031	1003.55734868488\\
160.860860860861	1015.68022257066\\
161.411411411411	1027.6872193994\\
161.961961961962	1039.57390508069\\
162.512512512513	1051.33588995412\\
163.063063063063	1062.9688304103\\
163.613613613614	1074.46843049494\\
164.164164164164	1085.83044349532\\
164.714714714715	1097.05067350856\\
165.265265265265	1108.12497699112\\
165.815815815816	1119.04926428899\\
166.366366366366	1129.81950114798\\
166.916916916917	1140.43171020352\\
167.467467467467	1150.88197244947\\
168.018018018018	1161.16642868541\\
168.568568568569	1171.28128094177\\
169.119119119119	1181.2227938824\\
169.66966966967	1190.98729618402\\
170.22022022022	1200.571181892\\
170.770770770771	1209.97091175202\\
171.321321321321	1219.18301451705\\
171.871871871872	1228.20408822932\\
172.422422422422	1237.03080147658\\
172.972972972973	1245.65989462238\\
173.523523523524	1254.08818100987\\
174.074074074074	1262.31254813852\\
174.624624624625	1270.32995881364\\
175.175175175175	1278.13745226791\\
175.725725725726	1285.73214525483\\
176.276276276276	1293.11123311345\\
176.826826826827	1300.27199080409\\
177.377377377377	1307.21177391474\\
177.927927927928	1313.92801963755\\
178.478478478478	1320.41824771529\\
179.029029029029	1326.68006135731\\
179.57957957958	1332.7111481246\\
180.13013013013	1338.50928078381\\
180.680680680681	1344.07231812971\\
181.231231231231	1349.39820577594\\
181.781781781782	1354.4849769137\\
182.332332332332	1359.33075303803\\
182.882882882883	1363.93374464155\\
183.433433433433	1368.29225187533\\
183.983983983984	1372.40466517659\\
184.534534534535	1376.26946586315\\
185.085085085085	1379.88522669421\\
185.635635635636	1383.25061239748\\
186.186186186186	1386.36438016222\\
186.736736736737	1389.22538009826\\
187.287287287287	1391.83255566062\\
187.837837837838	1394.18494403966\\
188.388388388388	1396.28167651668\\
188.938938938939	1398.12197878472\\
189.48948948949	1399.70517123448\\
190.04004004004	1401.03066920533\\
190.590590590591	1402.09798320119\\
191.141141141141	1402.90671907131\\
191.691691691692	1403.45657815585\\
192.242242242242	1403.74735739612\\
192.792792792793	1403.77894940961\\
193.343343343343	1403.55134252965\\
193.893893893894	1403.06462080968\\
194.444444444444	1402.31896399223\\
194.994994994995	1401.31464744257\\
195.545545545546	1400.05204204696\\
196.096096096096	1398.53161407576\\
196.646646646647	1396.75392501115\\
197.197197197197	1394.71963133987\\
197.747747747748	1392.42948431072\\
198.298298298298	1389.88432965716\\
198.848848848849	1387.08510728497\\
199.399399399399	1384.03285092516\\
199.94994994995	1380.72868775224\\
200.5005005005	1377.17383796793\\
201.051051051051	1373.36961435054\\
201.601601601602	1369.31742177024\\
202.152152152152	1365.01875667016\\
202.702702702703	1360.47520651386\\
203.253253253253	1355.68844919899\\
203.803803803804	1350.66025243776\\
204.354354354354	1345.39247310405\\
204.904904904905	1339.88705654772\\
205.455455455455	1334.14603587624\\
206.006006006006	1328.1715312038\\
206.556556556557	1321.96574886844\\
207.107107107107	1315.53098061722\\
207.657657657658	1308.86960275995\\
208.208208208208	1301.98407529155\\
208.758758758759	1294.8769409837\\
209.309309309309	1287.55082444572\\
209.85985985986	1280.00843115537\\
210.41041041041	1272.25254645974\\
210.960960960961	1264.28603454663\\
211.511511511512	1256.11183738683\\
212.062062062062	1247.73297364766\\
212.612612612613	1239.15253757823\\
213.163163163163	1230.37369786674\\
213.713713713714	1221.3996964703\\
214.264264264264	1212.23384741772\\
214.814814814815	1202.87953558565\\
215.365365365365	1193.34021544857\\
215.915915915916	1183.6194098031\\
216.466466466466	1173.72070846703\\
217.017017017017	1163.64776695365\\
217.567567567568	1153.4043051218\\
218.118118118118	1142.99410580212\\
218.668668668669	1132.42101340013\\
219.219219219219	1121.68893247648\\
219.76976976977	1110.80182630506\\
220.32032032032	1099.76371540935\\
220.870870870871	1088.57867607772\\
221.421421421421	1077.25083885806\\
221.971971971972	1065.7843870324\\
222.522522522523	1054.18355507209\\
223.073073073073	1042.45262707398\\
223.623623623624	1030.5959351784\\
224.174174174174	1018.61785796931\\
224.724724724725	1006.52281885729\\
225.275275275275	994.315284446073\\
225.825825825826	981.99976288301\\
226.376376376376	969.580802194279\\
226.926926926927	957.062988605307\\
227.477477477477	944.450944847137\\
228.028028028028	931.749328449279\\
228.578578578579	918.962830019717\\
229.129129129129	906.096171512711\\
229.67967967968	893.154104485013\\
230.23023023023	880.141408341152\\
230.780780780781	867.062888568441\\
231.331331331331	853.923374962347\\
231.881881881882	840.727719842882\\
232.432432432432	827.480796262685\\
232.982982982983	814.187496207436\\
233.533533533534	800.852728789285\\
234.084084084084	787.481418433968\\
234.634634634635	774.078503062212\\
235.185185185185	760.648932266262\\
235.735735735736	747.197665481991\\
236.286286286286	733.729670157433\\
236.836836836837	720.249919918341\\
237.387387387387	706.763392731467\\
237.937937937938	693.275069066236\\
238.488488488488	679.789930055501\\
239.039039039039	666.312955656042\\
239.58958958959	652.849122809517\\
240.14014014014	639.403403604508\\
240.690690690691	625.980763440371\\
241.241241241241	612.586159193561\\
241.791791791792	599.224537387059\\
242.342342342342	585.900832363716\\
242.892892892893	572.619964464002\\
243.443443443443	559.386838208969\\
243.993993993994	546.206340489053\\
244.544544544545	533.083338759382\\
245.095095095095	520.022679242262\\
245.645645645646	507.029185137506\\
246.196196196196	494.107654841263\\
246.746746746747	481.262860174008\\
247.297297297297	468.499544618357\\
247.847847847848	455.822421567285\\
248.398398398398	443.236172583615\\
248.948948948949	430.745445671016\\
249.4994994995	418.35485355764\\
250.05005005005	406.068971992627\\
250.600600600601	393.892338056329\\
251.151151151151	381.829448484798\\
251.701701701702	369.884758009179\\
252.252252252252	358.062677710612\\
252.802802802803	346.367573391253\\
253.353353353353	334.803763962019\\
253.903903903904	323.375519847645\\
254.454454454454	312.087061409662\\
255.005005005005	300.942557387806\\
255.555555555556	289.946123360625\\
256.106106106106	279.101820225515\\
256.656656656657	268.413652699151\\
257.207207207207	257.885567838538\\
257.757757757758	247.521453583395\\
258.308308308308	237.325137320379\\
258.858858858859	227.300384469651\\
259.409409409409	217.450897094345\\
259.95995995996	207.780312533417\\
260.510510510511	198.292202058408\\
261.061061061061	188.990069554601\\
261.611611611612	179.877350227063\\
262.162162162162	170.957409332051\\
262.712712712713	162.233540934255\\
263.263263263263	153.708966690298\\
263.813813813814	145.386834659048\\
264.364364364364	137.270218139029\\
264.914914914915	129.362114533494\\
265.465465465465	121.665444243495\\
266.016016016016	114.183049589413\\
266.566566566567	106.917693761301\\
267.117117117117	99.8720597984686\\
267.667667667668	93.0487495986498\\
268.218218218218	86.4502829571457\\
268.768768768769	80.0790966362827\\
269.319319319319	73.9375434655356\\
269.86986986987	68.0278914726498\\
270.42042042042	62.3523230460596\\
270.970970970971	56.9129341289746\\
271.521521521522	51.7117334453475\\
272.072072072072	46.7506417580716\\
272.622622622623	42.0314911596584\\
273.173173173173	37.5560243956604\\
273.723723723724	33.325894221088\\
274.274274274274	29.3426627900606\\
274.824824824825	25.6078010789151\\
275.375375375375	22.1226883429855\\
275.925925925926	18.8886116072574\\
276.476476476476	15.9067651910679\\
277.027027027027	13.1782502670777\\
277.577577577578	10.7040744545859\\
278.128128128128	8.48515144744747\\
278.678678678679	6.52230067664291\\
279.229229229229	4.81624700767059\\
279.77977977978	3.36762047285973\\
280.33033033033	2.17695603870438\\
280.880880880881	1.24469340830439\\
281.431431431431	0.571176858986554\\
281.981981981982	0.156655115165607\\
282.532532532533	0.00128125649228671\\
283.083083083083	0.105112661322015\\
283.633633633634	0.468110985526387\\
284.184184184184	1.09014217665108\\
284.734734734735	1.97097652342412\\
285.285285285285	3.11028874058288\\
285.835835835836	4.50765808900117\\
286.386386386386	6.16256853106458\\
286.936936936937	8.07440892123916\\
287.487487487487	10.2424732317627\\
288.038038038038	12.6659608133752\\
288.588588588589	15.3439766909925\\
289.139139139139	18.2755318942138\\
289.68968968969	21.4595438225401\\
290.24024024024	24.8948366451709\\
290.790790790791	28.580141735226\\
291.341341341341	32.5140981382509\\
291.891891891892	36.69525307479\\
292.442442442442	41.1220624768978\\
292.992992992993	45.7928915583486\\
293.543543543544	50.7060154183496\\
294.094094094094	55.8596196785344\\
294.644644644645	61.2518011529979\\
295.195195195195	66.8805685511284\\
295.745745745746	72.7438432129758\\
296.296296296296	78.8394598768842\\
296.846846846847	85.1651674791064\\
297.397397397397	91.7186299851045\\
297.947947947948	98.4974272522284\\
298.498498498499	105.499055923451\\
299.049049049049	112.72093035186\\
299.5995995996	120.160383555484\\
300.15015015015	127.814668202211\\
300.700700700701	135.68095762435\\
301.251251251251	143.756346862499\\
301.801801801802	152.037853738322\\
302.352352352352	160.522419955845\\
302.902902902903	169.206912230862\\
303.453453453453	178.08812344803\\
304.004004004004	187.162773845235\\
304.554554554555	196.427512224772\\
305.105105105105	205.878917190955\\
305.655655655656	215.51349841353\\
306.206206206206	225.32769791673\\
306.756756756757	235.317891393127\\
307.307307307307	245.480389542097\\
307.857857857858	255.811439432243\\
308.408408408408	266.307225887323\\
308.958958958959	276.963872895165\\
309.509509509509	287.777445039046\\
310.06006006006	298.743948951016\\
310.610610610611	309.859334786613\\
311.161161161161	321.11949772044\\
311.711711711712	332.520279462038\\
312.262262262262	344.05746979155\\
312.812812812813	355.726808114433\\
313.363363363363	367.523985034976\\
313.913913913914	379.44464394764\\
314.464464464464	391.484382645953\\
315.015015015015	403.638754948208\\
315.565565565566	415.903272339399\\
316.116116116116	428.273405628799\\
316.666666666667	440.744586622552\\
317.217217217217	453.31220981067\\
317.767767767768	465.971634067818\\
318.318318318318	478.718184367236\\
318.868868868869	491.547153507191\\
319.419419419419	504.453803849339\\
319.96996996997	517.433369068248\\
320.520520520521	530.481055911595\\
321.071071071071	543.592045970277\\
321.621621621622	556.761497457803\\
322.172172172172	569.984546998333\\
322.722722722723	583.256311422684\\
323.273273273273	596.571889571646\\
323.823823823824	609.926364105939\\
324.374374374374	623.314803322156\\
324.924924924925	636.732262973993\\
325.475475475475	650.173788098128\\
326.026026026026	663.634414844046\\
326.576576576577	677.109172307152\\
327.127127127127	690.593084364479\\
327.677677677678	704.081171512367\\
328.228228228228	717.568452705299\\
328.778778778779	731.049947195401\\
329.329329329329	744.520676371784\\
329.87987987988	757.975665599103\\
330.43043043043	771.409946054656\\
330.980980980981	784.818556563329\\
331.531531531532	798.196545429716\\
332.082082082082	811.53897226675\\
332.632632632633	824.840909820141\\
333.183183183183	838.097445787968\\
333.733733733734	851.303684634799\\
334.284284284284	864.45474939947\\
334.834834834835	877.545783496245\\
335.385385385385	890.571952508204\\
335.935935935936	903.528445972609\\
336.486486486487	916.410479157361\\
337.037037037037	929.213294827958\\
337.587587587588	941.932165004312\\
338.138138138138	954.562392706751\\
338.688688688689	967.099313690578\\
339.239239239239	979.53829816854\\
339.78978978979	991.874752520573\\
340.34034034034	1004.10412099018\\
340.890890890891	1016.22188736689\\
341.441441441441	1028.22357665393\\
341.991991991992	1040.10475672099\\
342.542542542543	1051.8610399408\\
343.093093093093	1063.48808480957\\
343.643643643644	1074.9815975502\\
344.194194194194	1086.33733369797\\
344.744744744745	1097.551099668\\
345.295295295295	1108.61875430385\\
345.845845845846	1119.53621040691\\
346.396396396396	1130.29943624568\\
346.946946946947	1140.90445704471\\
347.497497497497	1151.34735645242\\
348.048048048048	1161.62427798744\\
348.598598598599	1171.73142646269\\
349.149149149149	1181.66506938698\\
349.6996996997	1191.42153834337\\
350.25025025025	1200.99723034386\\
350.800800800801	1210.38860916002\\
351.351351351351	1219.5922066288\\
351.901901901902	1228.60462393335\\
352.452452452452	1237.42253285818\\
353.003003003003	1246.04267701821\\
353.553553553554	1254.46187306135\\
354.104104104104	1262.67701184411\\
354.654654654655	1270.68505957972\\
355.205205205205	1278.48305895854\\
355.755755755756	1286.06813024014\\
356.306306306306	1293.43747231681\\
356.856856856857	1300.58836374791\\
357.407407407407	1307.51816376495\\
357.957957957958	1314.22431324676\\
358.508508508508	1320.7043356646\\
359.059059059059	1326.95583799665\\
359.60960960961	1332.97651161181\\
360.16016016016	1338.76413312221\\
360.710710710711	1344.31656520433\\
361.261261261261	1349.63175738824\\
361.811811811812	1354.70774681488\\
362.362362362362	1359.5426589609\\
362.912912912913	1364.13470833091\\
363.463463463463	1368.48219911685\\
364.014014014014	1372.58352582425\\
364.564564564565	1376.43717386509\\
365.115115115115	1380.04172011718\\
365.665665665666	1383.39583344967\\
366.216216216216	1386.49827521463\\
366.766766766767	1389.34789970449\\
367.317317317317	1391.94365457513\\
367.867867867868	1394.28458123449\\
368.418418418418	1396.36981519659\\
368.968968968969	1398.19858640078\\
369.51951951952	1399.7702194961\\
370.07007007007	1401.08413409069\\
370.620620620621	1402.13984496612\\
371.171171171171	1402.9369622566\\
371.721721721722	1403.47519159291\\
372.272272272272	1403.75433421117\\
372.822822822823	1403.77428702616\\
373.373373373373	1403.53504266948\\
373.923923923924	1403.03668949219\\
374.474474474474	1402.27941153226\\
375.025025025025	1401.26348844652\\
375.575575575576	1399.98929540746\\
376.126126126126	1398.45730296464\\
376.676676676677	1396.66807687094\\
377.227227227227	1394.62227787361\\
377.777777777778	1392.32066147029\\
378.328328328328	1389.76407762998\\
378.878878878879	1386.95347047918\\
379.429429429429	1383.88987795322\\
379.97997997998	1380.57443141294\\
380.530530530531	1377.00835522693\\
381.081081081081	1373.19296631934\\
381.631631631632	1369.12967368358\\
382.182182182182	1364.81997786195\\
382.732732732733	1360.26547039155\\
383.283283283283	1355.46783321651\\
383.833833833834	1350.42883806687\\
384.384384384384	1345.15034580426\\
384.934934934935	1339.63430573478\\
385.485485485485	1333.88275488904\\
386.036036036036	1327.89781726996\\
386.586586586587	1321.68170306838\\
387.137137137137	1315.23670784686\\
387.687687687688	1308.56521169191\\
388.238238238238	1301.66967833509\\
388.788788788789	1294.55265424316\\
389.339339339339	1287.21676767767\\
389.88988988989	1279.66472772438\\
390.44044044044	1271.8993232928\\
390.990990990991	1263.92342208631\\
391.541541541542	1255.73996954309\\
392.092092092092	1247.35198774839\\
392.642642642643	1238.76257431856\\
393.193193193193	1229.97490125705\\
393.743743743744	1220.99221378309\\
394.294294294294	1211.81782913316\\
394.844844844845	1202.45513533605\\
395.395395395395	1192.90758996164\\
395.945945945946	1183.17871884407\\
396.496496496496	1173.27211477963\\
397.047047047047	1163.19143620004\\
397.597597597598	1152.94040582137\\
398.148148148148	1142.52280926926\\
398.698698698699	1131.94249368099\\
399.249249249249	1121.20336628466\\
399.7997997998	1110.30939295636\\
400.35035035035	1099.26459675558\\
400.900900900901	1088.07305643951\\
401.451451451451	1076.7389049568\\
402.002002002002	1065.2663279213\\
402.552552552553	1053.65956206633\\
403.103103103103	1041.92289368011\\
403.653653653654	1030.06065702283\\
404.204204204204	1018.07723272608\\
404.754754754755	1005.97704617512\\
405.305305305305	993.764565874531\\
405.855855855856	981.444301798198\\
406.406406406406	969.020803723635\\
406.956956956957	956.498659551907\\
407.507507507508	943.882493613305\\
408.058058058058	931.176964959627\\
408.608608608609	918.386765643623\\
409.159159159159	905.516618986261\\
409.70970970971	892.57127783244\\
410.26026026026	879.555522795799\\
410.810810810811	866.474160493275\\
411.361361361361	853.33202177006\\
411.911911911912	840.133959915612\\
412.462462462462	826.884848871337\\
413.013013013013	813.589581430738\\
413.563563563564	800.253067432507\\
414.114114114114	786.880231947367\\
414.664664664665	773.476013459285\\
415.215215215215	760.045362041725\\
415.765765765766	746.593237529626\\
416.316316316316	733.124607687774\\
416.866866866867	719.644446376245\\
417.417417417417	706.157731713596\\
417.967967967968	692.669444238485\\
418.518518518519	679.184565070401\\
419.069069069069	665.708074070167\\
419.61961961962	652.244948000924\\
420.17017017017	638.800158690257\\
420.720720720721	625.378671194098\\
421.271271271271	611.985441963234\\
421.821821821822	598.62541701288\\
422.372372372372	585.303530096173\\
422.922922922923	572.024700882169\\
423.473473473473	558.793833139049\\
424.024024024024	545.615812923199\\
424.574574574575	532.495506774823\\
425.125125125125	519.437759920771\\
425.675675675676	506.447394485232\\
426.226226226226	493.52920770897\\
426.776776776777	480.687970177698\\
427.327327327327	467.928424060436\\
427.877877877878	455.255281358136\\
428.428428428428	442.673222163679\\
428.978978978979	430.186892933507\\
429.52952952953	417.800904771733\\
430.08008008008	405.519831727294\\
430.630630630631	393.348209104791\\
431.181181181181	381.290531789638\\
431.731731731732	369.351252588134\\
432.282282282282	357.534780583078\\
432.832832832833	345.845479505527\\
433.383383383383	334.287666123313\\
433.933933933934	322.865608646855\\
434.484484484484	311.58352515303\\
435.035035035035	300.445582027357\\
435.585585585586	289.455892425471\\
436.136136136136	278.61851475413\\
436.686686686687	267.937451172481\\
437.237237237237	257.416646114096\\
437.787787787788	247.059984830321\\
438.338338338338	236.871291955489\\
438.888888888889	226.854330094503\\
439.439439439439	217.012798433344\\
439.98998998999	207.350331372985\\
440.540540540541	197.870497187246\\
441.091091091091	188.576796705026\\
441.641641641642	179.472662017518\\
442.192192192192	170.561455210739\\
442.742742742743	161.846467123942\\
443.293293293293	153.330916134334\\
443.843843843844	145.017946968557\\
444.394394394394	136.910629541363\\
444.944944944945	129.011957821919\\
445.495495495496	121.324848728158\\
446.046046046046	113.852141049583\\
446.596596596597	106.596594398926\\
447.147147147147	99.5608881930403\\
447.697697697698	92.7476206634182\\
448.248248248248	86.1593078966807\\
448.798798798799	79.798382905412\\
449.349349349349	73.6671947296483\\
449.8998998999	67.7680075694237\\
450.45045045045	62.1029999486\\
451.001001001001	56.6742639103592\\
451.551551551552	51.4838042446286\\
452.102102102102	46.5335377477274\\
452.652652652653	41.8252925145087\\
453.203203203203	37.3608072632583\\
453.753753753754	33.1417306935998\\
454.304304304304	29.1696208776424\\
454.854854854855	25.4459446845996\\
455.405405405405	21.972077239073\\
455.955955955956	18.7493014132579\\
456.506506506507	15.7788073531529\\
457.057057057057	13.0616920390756\\
457.607607607608	10.5989588805445\\
458.158158158158	8.39151734572934\\
458.708708708709	6.44018262559237\\
459.259259259259	4.74567533284416\\
459.80980980981	3.30862123582716\\
460.36036036036	2.12955102742426\\
460.910910910911	1.20890012907811\\
461.461461461461	0.547008529993268\\
462.012012012012	0.144120661580943\\
462.562562562563	0.000385307192305361\\
463.113113113113	0.115855547173805\\
463.663663663664	0.490488739265538\\
464.214214214214	1.12414653434819\\
464.764764764765	2.01659492753444\\
465.315315315315	3.16750434458515\\
465.865865865866	4.57644976361856\\
466.416416416416	6.24291087206757\\
466.966966966967	8.1662722588271\\
467.517517517517	10.3458236415206\\
468.068068068068	12.780760128802\\
468.618618618619	15.4701825175953\\
469.169169169169	18.413097625164\\
469.71971971972	21.6084186558849\\
470.27027027027	25.0549656025916\\
470.820820820821	28.7514656823513\\
471.371371371371	32.696553806479\\
471.921921921922	36.8887730846642\\
472.472472472472	41.3265753629861\\
473.023023023023	46.0083217956352\\
473.573573573574	50.9322834501263\\
474.124124124124	56.0966419457797\\
474.674674674675	61.4994901252349\\
475.225225225225	67.1388327587494\\
475.775775775776	73.0125872810207\\
476.326326326326	79.1185845602578\\
476.876876876877	85.4545696992468\\
477.427427427427	92.0182028680167\\
477.977977977978	98.8070601679854\\
478.528528528529	105.818634527035\\
479.079079079079	113.050336625376\\
479.62962962963	120.499495851762\\
480.18018018018	128.163361289722\\
480.730730730731	136.039102733453\\
481.281281281281	144.123811732995\\
481.831831831832	152.414502668297\\
482.382382382382	160.908113851784\\
482.932932932933	169.601508659013\\
483.483483483483	178.491476687\\
484.034034034034	187.574734939822\\
484.584584584585	196.84792904096\\
485.135135135135	206.307634472074\\
485.685685685686	215.950357837643\\
486.236236236236	225.772538155048\\
486.786786786787	235.770548169615\\
487.337337337337	245.940695694134\\
487.887887887888	256.279224972353\\
488.438438438438	266.78231806595\\
488.988988988989	277.446096264468\\
489.53953953954	288.266621517692\\
490.09009009009	299.239897889943\\
490.640640640641	310.361873035733\\
491.191191191191	321.6284396963\\
491.741741741742	333.035437216339\\
492.292292292292	344.578653080531\\
492.842842842843	356.253824469182\\
493.393393393393	368.056639832451\\
493.943943943944	379.982740482573\\
494.494494494494	392.027722203486\\
495.045045045045	404.187136877275\\
495.595595595596	416.456494126826\\
496.146146146146	428.831262974089\\
496.696696696697	441.306873513334\\
497.247247247247	453.878718598776\\
497.797797797798	466.542155545955\\
498.348348348348	479.292507846287\\
498.898898898899	492.125066893958\\
499.44944944945	505.035093724922\\
500	518.017820766804\\
};
\addlegendentry{Fit}

\addplot[only marks, mark=*, mark options={}, mark size=1.5000pt, draw=mycolor2, fill=red] table[row sep=crcr]{%
x	y\\
281	2.9\\
90	91\\
45	1000\\
92	19\\
};
\addlegendentry{Data}

\end{axis}
\end{tikzpicture}%
\end{figure}
The figure is fit using the equation
\[y =  k\cos\left(\frac{\pi}{180}(x+c)\right)^{2}\]
and we get the parameters as \(k=1403.79748\) and \(c=167.42097\).
The data is not centered around zero as the laser source can be oriented in any angle which allows $c$ to be free and $k$ is just the amplitude.
Thus, the detector does show the zero reading at $102.57903^\circ$ but due to the least count of the apparatus, we cannot reach those precise values.
We now proceed ot calculate the value for brewster angle.
\begin{figure}[H]
    \centering
    \setlength\fwidth{0.75\textwidth}
    \setlength\fheight{0.5\textwidth}
    % This file was created by matlab2tikz.
%
\definecolor{mycolor1}{rgb}{0.00000,0.44700,0.74100}%
\definecolor{mycolor2}{rgb}{0.12941,0.12941,0.12941}%
%
\begin{tikzpicture}

\begin{axis}[%
width=0.951\fwidth,
height=\fheight,
at={(0\fwidth,0\fheight)},
scale only axis,
xmin=0,
xmax=5,
xtick={0, 1, 2, 3, 4, 5},
xlabel style={font=\color{mycolor2}},
xlabel={Measurement Number},
ymin=332,
ymax=348,
ylabel style={font=\color{mycolor2}},
ylabel={Value},
axis background/.style={fill=white},
xmajorgrids,
ymajorgrids,
legend style={legend cell align=left, align=left}
]
\addplot [color=mycolor1, only marks, mark=o, mark options={solid, mycolor1}]
 plot [error bars/.cd, y dir=both, y explicit, error bar style={line width=0.5pt}, error mark options={line width=0.5pt, mark size=6.0pt, rotate=90}]
 table[row sep=crcr, y error plus index=2, y error minus index=3]{%
1	343	1	1\\
2	342	1	1\\
3	339	4	4\\
4	343	1	1\\
};
\addlegendentry{Data}

\addplot [color=red]
  table[row sep=crcr]{%
0	342.591836734694\\
6	342.591836734694\\
};
\addlegendentry{Weighted Average}

\end{axis}
\end{tikzpicture}%
    \caption{The value is as $ 342.59 \pm 0.57$ }
\end{figure}
We subtract the value obtained above from $39^\circ$ to get the value for brewster angle.
\section*{Result}
We obtain the value of brewster angle as $ 56.41 \pm 0.57$ degrees.
\end{document}